%% toc: P2 10 10.5
\section{Connective idioms, and what first-order sentences can say}
\label{sec:p02-w02-p02-005-idioms}

\Cref{sec:prop-syntax-semantics-03} fixes the method of translation from
controlled prose: a key assigns one propositional variable to each atomic
statement, and the translation of a sentence is a formula together with that
key.  Its table treats ``if'', ``only if'', ``iff'' and ``unless'' placed after
the main clause, and it treats exclusive disjunction.  This section supplies
the remaining rows of the translation guide for the idioms ``only'' and
``neither \dots\ nor'', proves that the rows are correct and that the nearby
mistranslations are not equivalent to them, and then turns to the question
one level up: which mathematical notions a first-order sentence, or a set of
first-order sentences, can express at all.  For that question the section
follows the expository chapter of Barwise \cite[ch.~A.1, pp.~5--46]{Barwise1977FirstOrder}.

\emph{Inputs.}  Propositional formulas, valuations and the key-based
translation method of \cref{sec:prop-syntax-semantics-03}; for the
first-order part, the compactness theorem \cref{thm:fol-compactness} and the
downward L\"owenheim--Skolem theorem \cref{thm:downward-ls}.
\emph{Conclusions.}  Canonical translations of the fronted ``only if'',
of ``neither \dots\ nor'' and of ``not both'', with proofs that each is the
correct one (\cref{thm:p02-w02-p02-005-idiom-equivalences}); the precise
syntax and semantics of first-order logic in the form Barwise gives them;
and his results that divisible, torsion-free and torsion abelian groups,
the ordered field of real numbers, fields of characteristic \(0\) and
principal ideal rings escape finite or arbitrary first-order axiomatization,
while maximality of an ideal does not; then Barwise's reduction of
first-order logic to propositional logic by Henkin constants, his Hilbert
and Gentzen systems with their completeness theorems
(\cref{thm:p02-w02-p02-005-bw-completeness-h,thm:p02-w02-p02-005-bw-completeness-g}),
and the compactness and L\"owenheim--Skolem behaviour of the extensions of
first-order logic in his \S5.  Two places where his \S4 needs repair are
marked and repaired: Henkin axioms for vacuous quantifiers
(\cref{thm:p02-w02-p02-005-bw-vacuous}), and Henkin axioms that share a
witnessing constant (\cref{thm:p02-w02-p02-005-bw-lemma-411}).
\emph{Ambient assumptions.}  The metatheory is ZFC; choice enters through
\cref{thm:fol-compactness,thm:downward-ls}, through the choice of witnesses
in \cref{thm:p02-w02-p02-005-bw-lemma-47} and through the well-ordering in
\cref{thm:p02-w02-p02-005-bw-prop-compactness}.  Structures have nonempty
domains.

\subsection{The idioms ``only'' and ``neither \dots\ nor''}
\label{sec:p02-w02-p02-005-only-neither}

\begin{hbconvention}{def:p02-w02-p02-005-only}{The scope of ``only'' in the propositional guide}
In the controlled fragment translated with a propositional key, the word
``only'' is a sentential connective exactly when it forms the compound
``only if'' joining two clauses.  The compound may follow the main clause
(``\(A\) only if \(B\)'') or stand in front of the sentence
(``Only if \(B\), \(A\)'', and its inverted variant ``Only if \(B\) is
\(A\)'').  In each position the clause introduced by ``only if'' is a
\emph{necessary} condition for the other clause, and the translation is
\(A\to B\), where \(A\) and \(B\) are the translations of the two clauses.
Moving the ``only if'' clause to the front changes the word order and
leaves the formula unchanged.

Every other occurrence of ``only'' modifies a noun phrase or a predicate,
as in ``Only valid archives are readable''.  Such an occurrence restricts
a quantity of objects; it is not a truth function of whole clauses, so a
propositional key cannot translate it.  With a first-order key \(V(x)\):
\(x\) is a valid archive, \(R(x)\): \(x\) is readable, the sentence
translates as \(\forall x\,(R(x)\to V(x))\): the restricted noun phrase
supplies the consequent, exactly as the ``only if'' clause does.
Translating it as \(\forall x\,(V(x)\to R(x))\) states that every valid
archive is readable, which is a different claim.
\index{only if!fronted}\index{translation!only}
\end{hbconvention}

\begin{hbdefinition}{def:p02-w02-p02-005-neither}{Neither \dots\ nor, and not both}
``Neither \(A\) nor \(B\)'' asserts that \(A\) is false and that \(B\) is
false.  Its canonical translation is \(\neg A\land\neg B\).  The formula
\(\neg(A\lor B)\) is an equivalent alternative
(\cref{thm:p02-w02-p02-005-idiom-equivalences}).  ``Not both \(A\) and
\(B\)'' denies only the conjunction, and its canonical translation is
\(\neg(A\land B)\), equivalently \(\neg A\lor\neg B\).
\index{neither nor@neither \dots\ nor}\index{not both}
\end{hbdefinition}

With the key of \cref{sec:prop-syntax-semantics-03}, \(P\): the archive is
readable, \(Q\): the checksum is valid, the guide gains the rows of \cref{tab:p02-w02-p02-005-idioms}.
The ``unless'' row is repeated with its simplified form so that the two
readings of ``unless'' can be compared; the formula in the middle column is
the translation, and the right column records a classically equivalent
simplification, which is a further step.

\begin{table}[htbp]
\centering
\small
\caption{Translation-guide rows for front-position ``only if'', ``unless'',
``neither \dots\ nor'' and ``not both''.}
\label{tab:p02-w02-p02-005-idioms}
\begin{tabular}{@{}p{0.42\textwidth}p{0.2\textwidth}p{0.3\textwidth}@{}}
\toprule
Controlled prose & Translation & Equivalent form, and the error it excludes \\
\midrule
Only if the checksum is valid is the archive readable. & \(P\to Q\) &
\(\neg Q\to\neg P\); not \(Q\to P\) \\
Only if the checksum is valid, the archive is readable. & \(P\to Q\) &
\(\neg Q\to\neg P\); not \(Q\to P\) \\
The archive is readable unless the checksum is invalid. & \(\neg\neg Q\to P\) &
\(Q\to P\), and \(\neg Q\lor P\) \\
Neither is the archive readable nor is the checksum valid. &
\(\neg P\land\neg Q\) & \(\neg(P\lor Q)\); not \(\neg P\lor\neg Q\) \\
The archive is not both readable and accompanied by a valid checksum. &
\(\neg(P\land Q)\) & \(\neg P\lor\neg Q\); strictly weaker than the
neither row \\
\bottomrule
\end{tabular}
\end{table}

\begin{hbtheorem}{thm:p02-w02-p02-005-idiom-equivalences}{proved=annex:p02-w02-p02-005-idiom-equivalences}{Correctness of the idiom rows}
Let \(A\) and \(B\) be propositional formulas, and write
\(\varphi\equiv\psi\) when every valuation gives \(\varphi\) and \(\psi\)
the same value.
\begin{enumerate}
\item \(A\to B\equiv\neg B\to\neg A\).  If some valuation makes \(A\) false
      and \(B\) true, then \(A\to B\not\equiv B\to A\).
\item \(\neg\neg B\to A\equiv B\to A\equiv\neg B\lor A\).
\item \(\neg A\land\neg B\equiv\neg(A\lor B)\).
\item Every valuation making \(\neg A\land\neg B\) true makes
      \(\neg(A\land B)\) true.  If some valuation makes exactly one of
      \(A\) and \(B\) true, then \(\neg(A\land B)\not\equiv\neg A\land\neg B\),
      and \(\neg A\lor\neg B\not\equiv\neg A\land\neg B\).
\end{enumerate}
\end{hbtheorem}

\begin{hbexample}{def:p02-w02-p02-005-idiom-example}{Checking the rows against a valuation}
Take the valuation with \(v(P)=1\) and \(v(Q)=0\): the archive is readable
and the checksum is invalid.  ``Only if the checksum is valid is the archive
readable'' is then false, and \(P\to Q\) is false, as it should be.  The
mistranslation \(Q\to P\) is true at this valuation.  ``Neither is the
archive readable nor is the checksum valid'' is false, since the archive is
readable; \(\neg P\land\neg Q\) is false, while \(\neg P\lor\neg Q\) is true.
At this valuation ``not both'' is true and ``neither'' is false, which is
the separation that clause 4 of \cref{thm:p02-w02-p02-005-idiom-equivalences}
records.
\end{hbexample}

\begin{hbwarning}{def:p02-w02-p02-005-idiom-limits}{Where the rows stop}
The equivalences in the right-hand column are equivalences of classical
propositional logic.  The step from \(\neg\neg Q\to P\) to \(Q\to P\)
uses double negation and fails intuitionistically, which is why the
translation keeps the double negation.  The rows translate truth-functional
uses only.  ``Only'' attached to a noun phrase needs a first-order key
(\cref{def:p02-w02-p02-005-only}).  The temporal or causal readings that
``only if'' and ``unless'' can carry in ordinary prose are excluded from
the controlled fragment.
\end{hbwarning}

\subsection{First-order languages, structures and satisfaction after Barwise}
\label{sec:p02-w02-p02-005-barwise-setup}

Barwise's chapter is written for a mathematician who has not studied logic
\cite[ch.~A.1, \S1, p.~6]{Barwise1977FirstOrder}.  Its reasoning task is to
recognize, for a given mathematical notion, whether a first-order sentence
or set of sentences expresses it, and to prove the negative answers.  The
following definitions are those of its \S3, in its notation
\cite[ch.~A.1, \S3, pp.~17--22]{Barwise1977FirstOrder}.

\begin{hbdefinition}{def:p02-w02-p02-005-bw-language}{Language and logical symbols}
The logical symbols are the connectives \(\land\), \(\lor\), \(\neg\),
\(\to\), the equality symbol \(=\), the quantifiers \(\forall\) and
\(\exists\), an infinite sequence of variables \(x,y,z,x_1,y_1,\ldots\),
and the parentheses \()\), \((\) \cite[ch.~A.1, \S2, pp.~6--7]{Barwise1977FirstOrder}.
A \emph{language} \(L\) is a set of function symbols, relation symbols and
constant symbols, of any size.  Each function symbol \(f\in L\) carries a
positive integer \(\#(f)\); if \(n=\#(f)\), then \(f\) is \emph{\(n\)-ary}.
Each relation symbol \(R\in L\) likewise carries a positive integer
\(\#(R)\).  An \emph{expression} is a finite sequence each of whose entries
is a logical symbol or a member of \(L\)
\cite[ch.~A.1, \S3, pp.~17--18]{Barwise1977FirstOrder}.
For example, group theory uses \(L=\{+,0\}\) with \(\#(+)=2\), and set
theory uses \(L=\{\in\}\) with \(\#(\in)=2\).
\end{hbdefinition}

\begin{hbdefinition}{def:p02-w02-p02-005-bw-structure}{Structure for a language (Barwise, Definition~3.1)}
A (set-theoretic) \emph{structure for \(L\)} is a pair
\(\mathfrak M=(M,F)\) in which \(M\) is a nonempty set and \(F\) is a
function with domain \(L\).  Writing \(x^{\mathfrak M}\) for \(F(x)\):
if \(R\in L\) is an \(n\)-ary relation symbol, then
\(R^{\mathfrak M}\subseteq M^n\); if \(f\in L\) is an \(n\)-ary function
symbol, then \(f^{\mathfrak M}:M^n\to M\); if \(c\in L\) is a constant
symbol, then \(c^{\mathfrak M}\in M\).  One writes
\(\mathfrak M=(M,R^{\mathfrak M},\ldots,f^{\mathfrak M},\ldots,c^{\mathfrak M},\ldots)\)
\cite[ch.~A.1, \S3, Definition~3.1, p.~17]{Barwise1977FirstOrder}.
The chapter considers only set-theoretic structures; a structure whose
domain is a proper class, such as the universe of all sets for
\(L=\{\in\}\), is set aside \cite[ch.~A.1, \S3, pp.~17--18]{Barwise1977FirstOrder}.
\end{hbdefinition}

\begin{hbdefinition}{def:p02-w02-p02-005-bw-syntax}{Terms, formulas and sentences (Barwise, Definitions~3.2--3.6)}
The \emph{terms} of \(L\) form the smallest set of expressions that contains
the variables and the constant symbols of \(L\) and contains
\(f(t_1\cdots t_n)\) whenever \(t_1,\ldots,t_n\) are terms and \(f\in L\) is
\(n\)-ary.  A \emph{closed term} contains no variable.  An \emph{atomic
formula} is an expression \((t_1=t_2)\), with terms \(t_1,t_2\), or
\(R(t_1\cdots t_n)\), with \(R\in L\) an \(n\)-ary relation symbol and terms
\(t_1,\ldots,t_n\).  The \emph{formulas} of \(L\) form the smallest set of
expressions that contains the atomic formulas and contains
\(\neg\varphi\), \((\varphi\land\psi)\), \((\varphi\lor\psi)\),
\((\varphi\to\psi)\) whenever it contains \(\varphi\) and \(\psi\), and
\((\exists v\varphi)\), \((\forall v\varphi)\) whenever it contains
\(\varphi\) and \(v\) is a variable.  Repeated binary symbols associate to
the right: \(\varphi\land\psi\land\theta\) is
\((\varphi\land(\psi\land\theta))\) and \(\varphi\to\psi\to\theta\) is
\((\varphi\to(\psi\to\theta))\).

The set \(\FV(\varphi)\) of \emph{free variables} is the set of variables
occurring in \(\varphi\) when \(\varphi\) is atomic;
\(\FV(\neg\varphi)=\FV(\varphi)\);
\(\FV(\varphi\land\psi)=\FV(\varphi\lor\psi)=\FV(\varphi\to\psi)
=\FV(\varphi)\cup\FV(\psi)\); and
\(\FV(\exists v\varphi)=\FV(\forall v\varphi)=\FV(\varphi)-\{v\}\).
Writing \(\varphi(v_1,\ldots,v_n)\) asserts only
\(\FV(\varphi)\subseteq\{v_1,\ldots,v_n\}\).  A \emph{sentence} is a
formula with no free variables
\cite[ch.~A.1, \S3, Definitions~3.2--3.6, pp.~18--20]{Barwise1977FirstOrder}.
\end{hbdefinition}

The abbreviations of the chapter are used below: \(x+y\) for \(+(xy)\),
\(nx\) for the term \((\cdots((x+x)+x)+\cdots+x)\) with \(n\geq1\)
occurrences of \(x\), and \(n1\) for \((\cdots((1+1)+1)+\cdots+1)\) with
\(n\) occurrences of \(1\) \cite[ch.~A.1, \S2, pp.~8, 12; \S3, p.~18]{Barwise1977FirstOrder}.

\begin{hbdefinition}{def:p02-w02-p02-005-bw-satisfaction}{Assignments, satisfaction and models (Barwise, Definitions~3.7--3.8)}
Let \(\mathfrak M=(M,\ldots)\) be a structure for \(L\).  An
\emph{assignment} in \(\mathfrak M\) is a function \(s\) from the variables
into \(M\), and \(s(\tfrac{a}{v})\) is the assignment that agrees with
\(s\) except that it sends \(v\) to \(a\).  The value
\(t^{\mathfrak M}(s)\) of a term is \(c^{\mathfrak M}\) for a constant
symbol \(c\), \(s(v)\) for a variable \(v\), and
\(f^{\mathfrak M}(t_1^{\mathfrak M}(s),\ldots,t_n^{\mathfrak M}(s))\) for
\(f(t_1\cdots t_n)\).  The relation \(\mathfrak M\models\varphi[s]\)
(the assignment \(s\) satisfies \(\varphi\) in \(\mathfrak M\)) is defined
by recursion on \(\varphi\):
\(\mathfrak M\models(t_1=t_2)[s]\) iff
\(t_1^{\mathfrak M}(s)=t_2^{\mathfrak M}(s)\);
\(\mathfrak M\models R(t_1\cdots t_n)[s]\) iff
\((t_1^{\mathfrak M}(s),\ldots,t_n^{\mathfrak M}(s))\in R^{\mathfrak M}\);
\(\mathfrak M\models\neg\varphi[s]\) iff not
\(\mathfrak M\models\varphi[s]\);
\(\mathfrak M\models(\varphi\land\psi)[s]\) iff both conjuncts are
satisfied by \(s\);
\(\mathfrak M\models(\varphi\lor\psi)[s]\) iff at least one disjunct is;
\(\mathfrak M\models(\varphi\to\psi)[s]\) iff \(\varphi\) is not satisfied
by \(s\) or \(\psi\) is;
\(\mathfrak M\models(\exists v\varphi)[s]\) iff
\(\mathfrak M\models\varphi[s(\tfrac{a}{v})]\) for some \(a\in M\);
and \(\mathfrak M\models(\forall v\varphi)[s]\) iff
\(\mathfrak M\models\varphi[s(\tfrac{a}{v})]\) for all \(a\in M\)
\cite[ch.~A.1, \S3, Definitions~3.7--3.8, pp.~20--21]{Barwise1977FirstOrder}.

For a sentence \(\varphi\), \(\mathfrak M\models\varphi\) means that some,
equivalently every, assignment satisfies \(\varphi\); \(\mathfrak M\) is
then a \emph{model} of \(\varphi\), and a model of a set \(\Phi\) of
sentences if it is a model of each member.  Structures \(\mathfrak M\) and
\(\mathfrak N\) for \(L\) are \emph{elementarily equivalent},
\(\mathfrak M\equiv\mathfrak N\), if they are models of the same sentences
of \(L\).  A class \(K\) of structures for \(L\) is \emph{axiomatizable}
(\emph{finitely axiomatizable}) if some set (finite set) \(\Phi\) of
sentences of \(L\) has exactly the members of \(K\) as its models.
For a set \(T\cup\{\psi\}\) of sentences, \(T\models\psi\) means that every
model of \(T\) is a model of \(\psi\)
\cite[ch.~A.1, \S2, p.~10; \S3, pp.~21--22]{Barwise1977FirstOrder}.
\end{hbdefinition}

\begin{hbconvention}{def:p02-w02-p02-005-bw-substitution}{Substitution with the bound-variable proviso}
For a formula \(\varphi(v)\) and a term \(t\), \(\varphi(t/v)\), or
\(\varphi(t)\), is the result of replacing every free occurrence of \(v\)
by \(t\).  The notation is used only when no variable of \(t\) occurs bound
in \(\varphi\); bound variables are renamed first otherwise.  Barwise's
example: for \(\varphi(v)=\exists w\,(v\neq w)\) and \(t=w\), the intended
\(\varphi(t)\) is \(\exists w'\,(w\neq w')\), whereas the unrenamed
\(\exists w\,(w\neq w)\) says something false in every structure
\cite[ch.~A.1, \S3, p.~22]{Barwise1977FirstOrder}.
\end{hbconvention}

Two facts are stated without proof in the chapter after
Definitions~3.7 and~3.8: the value \(t^{\mathfrak M}(s)\) depends only on
the values of \(s\) at the variables occurring in \(t\), and whether
\(\mathfrak M\models\varphi[s]\) depends only on the values of \(s\) at the
members of \(\FV(\varphi)\) \cite[ch.~A.1, \S3, pp.~21--22]{Barwise1977FirstOrder}.
The second fact is what makes \(\mathfrak M\models\varphi\) well defined
for sentences.  The chapter also notes that isomorphic structures are
elementarily equivalent \cite[ch.~A.1, \S3, p.~22]{Barwise1977FirstOrder}.

\subsection{Deciding whether a notion is first-order}
\label{sec:p02-w02-p02-005-barwise-realm}

In first-order logic the quantifiers range over the elements of the domain
of the structure.  Second-order logic also quantifies over subsets of the
domain and over functions on it; weak second-order logic quantifies over
finite subsets and natural numbers \cite[ch.~A.1, \S2, p.~7]{Barwise1977FirstOrder}.
A sentence that quantifies over the natural numbers, or over subsets, is
therefore not first-order, even though it may be equivalent on the
intended structures to an infinite list of first-order sentences.

\begin{hbexample}{def:p02-w02-p02-005-bw-groups}{Group-theoretic notions (Barwise, \S2)}
A group \(G=(G,+,0)\), with \(G\) nonempty, \(0\in G\) and
\(+:G\times G\to G\), is a model of
\[
\begin{aligned}
&\forall x\,\forall y\,\forall z\,[x+(y+z)=(x+y)+z], &&(1)\\
&\forall x\,[x+0=x], &&(2)\\
&\forall x\,\exists y\,[x+y=0]; &&(3)
\end{aligned}
\]
it is \emph{abelian} if it also models
\(\forall x\,\forall y\,[x+y=y+x]\) (4).  An abelian group has every element
of order \(\leq n\) if it models the single first-order sentence
\(\forall x\,[x=0\lor 2x=0\lor\cdots\lor nx=0]\) (5).  It is
\emph{divisible} if \(\forall n\geq1\,\forall x\,\exists y\,[ny=x]\) (6),
which quantifies over positive integers and so is not first-order; it is
replaced by the list of sentences \((6)_n\):
\(\forall x\,\exists y\,[ny=x]\) for \(n\geq2\), the case \(n=1\) being
trivially true.  It is \emph{torsion-free} if
\(\forall n\geq1\,\forall x\,[x\neq0\to nx\neq0]\) (7), replaced by the
list \((7)_n\): \(\forall x\,[x\neq0\to nx\neq0]\).  It is \emph{torsion}
if \(\forall x\,\exists n\geq1\,[nx=0]\) (8).  The attempted replacement
\(\forall x\,[x=0\lor 2x=0\lor\cdots\lor nx=0\lor\cdots]\) is an infinitely
long disjunction, which belongs to an infinitary logic rather than to
first-order logic \cite[ch.~A.1, \S2, pp.~7--9]{Barwise1977FirstOrder}.
Thus the notions of group, abelian group and abelian group of bounded
exponent are finitely axiomatized; divisible and torsion-free need infinite lists; torsion is
not first-order at all.  The three negative claims are
\cref{thm:p02-w02-p02-005-bw-divisible,thm:p02-w02-p02-005-bw-torsion-free,thm:p02-w02-p02-005-bw-torsion}.
\end{hbexample}

The chapter uses two tools for the negative results.  It proves both from
its reduction of first-order logic to propositional logic in its \S4; the
annex proofs give that route, through the Main Lemma
\cref{thm:p02-w02-p02-005-bw-main-lemma}, and also derive both from
\cref{thm:fol-compactness,thm:downward-ls}.

\begin{hbtheorem}{thm:p02-w02-p02-005-bw-compactness}{proved=annex:p02-w02-p02-005-bw-compactness}{Compactness, two forms (Barwise, 2.4)}
Let \(T\) be a set of first-order sentences of \(L\).
\begin{enumerate}
\item If every finite \(T_0\subseteq T\) has a model, then \(T\) has a model.
\item For every sentence \(\psi\) of \(L\), if \(T\models\psi\), then
      \(T_0\models\psi\) for some finite \(T_0\subseteq T\).
\end{enumerate}
Each form follows from the other by a one-step argument.
\end{hbtheorem}

\begin{hbtheorem}{thm:p02-w02-p02-005-bw-ls}{proved=annex:p02-w02-p02-005-bw-ls}{L\"owenheim--Skolem (Barwise, Theorem~2.5)}
Let \(\kappa\) be an infinite cardinal and let \(T\) be a set of at most
\(\kappa\) first-order sentences.  If \(T\) has a model, then \(T\) has a
model whose domain has cardinality at most \(\kappa\).  In particular, a
set of sentences of a finite or countable language that has a model has a
countable model.
\end{hbtheorem}

\begin{hbtheorem}{thm:p02-w02-p02-005-bw-divisible}{proved=annex:p02-w02-p02-005-bw-divisible}{Divisibility is not finitely axiomatizable (Barwise, Proposition~2.1)}
Any finite set of first-order sentences of \(\{+,0\}\) true in all
divisible abelian groups is true in some abelian group that is not
divisible.
\end{hbtheorem}

\begin{hbtheorem}{thm:p02-w02-p02-005-bw-torsion-free}{proved=annex:p02-w02-p02-005-bw-torsion-free}{Torsion-freeness is not finitely axiomatizable (Barwise, Proposition~2.2)}
The class of torsion-free abelian groups is not finitely axiomatizable in
first-order logic.
\end{hbtheorem}

\begin{hbtheorem}{thm:p02-w02-p02-005-bw-torsion}{proved=annex:p02-w02-p02-005-bw-torsion}{Torsion is not axiomatizable (Barwise, Proposition~2.3)}
The set of first-order sentences true in all torsion abelian groups is true
in some abelian group \(H\) that is not torsion.  More exactly, if
\(G=(G,+,0)\) is an abelian group such that for every \(n\) some element
of \(G\) has order \(\geq n\), then some abelian group \(H\) that is not
torsion satisfies \(G\equiv H\).  Hence the class of torsion abelian groups
is not axiomatizable by any set of first-order sentences, finite or
infinite.
\end{hbtheorem}

\begin{hbtheorem}{thm:p02-w02-p02-005-bw-reals}{proved=annex:p02-w02-p02-005-bw-reals}{No first-order categorical axioms for the reals (Barwise, Proposition~2.6)}
Let \(\mathbb R=(\mathbb R,+,\cdot,<,0,1)\) be the ordered field of real
numbers.  No set of first-order sentences of \(\{+,\cdot,<,0,1\}\) has as
its models exactly the structures isomorphic to \(\mathbb R\).
\end{hbtheorem}

\begin{hbtheorem}{thm:p02-w02-p02-005-bw-nonstandard}{proved=annex:p02-w02-p02-005-bw-nonstandard}{A non-Archimedean elementary extension (Barwise, Proposition~2.7)}
Expand \(\mathbb R\) by a constant symbol for every real number \(r\),
interpreted as \(r\).  Some non-Archimedean ordered field \({}^*\mathbb R\)
extending \(\mathbb R\) satisfies every first-order sentence of this
expanded language that is true in \(\mathbb R\).
\end{hbtheorem}

The chapter's diagnosis \cite[ch.~A.1, \S2, pp.~12--13]{Barwise1977FirstOrder}:
the Archimedean axiom \(\forall x\,\exists n\,[x\leq n1]\) quantifies over
the natural numbers, and the order-completeness axiom
\(\forall X\subseteq\mathbb R\,[\text{if }X\neq\emptyset\text{ is bounded,
then }X\text{ has a l.u.b.}]\) quantifies over all subsets of \(\mathbb R\).
Neither is first-order, and the characterization of \(\mathbb R\) up to
isomorphism holds only relative to a universe of sets.  In
\({}^*\mathbb R\) the element \(1/c\) is a positive infinitesimal, which is
the starting point of Robinson's nonstandard analysis; the chapter refers
to its Chapter~A.6 for this.

\begin{hbexample}{def:p02-w02-p02-005-bw-rings}{Rings, ideals and fields (Barwise, \S2)}
In the language \(\{+,\cdot,0,1\}\), a commutative ring with identity is a
model of (1)--(4) together with
\[
\begin{gathered}
\forall x\forall y\,[x\cdot y=y\cdot x],\qquad
\forall x\forall y\forall z\,[(x\cdot y)\cdot z=x\cdot(y\cdot z)],\\
\forall x\forall y\forall z\,[x\cdot(y+z)=(x\cdot y)+(x\cdot z)],\qquad
\forall x\,[x\cdot1=x],\qquad 0\neq1.
\end{gathered}
\]
An integral domain also models
\(\forall x\forall y\,[x\cdot y=0\to(x=0\lor y=0)]\).  To speak of an ideal,
add a unary relation symbol \(I\) and consider \((\mathfrak R,I)\).  The
subset \(I\) is a (proper) ideal exactly when \((\mathfrak R,I)\) models
\[
\begin{gathered}
I(0)\land\neg I(1),\qquad
\forall x\forall y\,[I(x)\land I(y)\to I(x+y)],\\
\forall x\forall y\,[I(x)\land x+y=0\to I(y)],\qquad
\forall x\forall y\,[I(y)\to I(x\cdot y)];
\end{gathered}
\]
it is prime if also \(\forall x\forall y\,[I(x\cdot y)\to I(x)\lor I(y)]\).
Maximality, \(\forall J\,[I\subseteq J\land J\text{ an ideal}\to J=I\lor
J=R]\) (9), quantifies over subsets, yet it has the first-order equivalent
\[
\forall x\,[\neg I(x)\to\exists y\,\exists z\,(I(z)\land x y+z=1)]
\tag*{(9)\('\)}
\]
(\cref{thm:p02-w02-p02-005-bw-maximal}).  Being a principal ideal ring,
\(\forall I\,[I\text{ an ideal}\to\exists x\,\forall y\,(I(y)\leftrightarrow
\exists z\,(y=zx))]\), has the same shape as (9) but no first-order
equivalent (\cref{thm:p02-w02-p02-005-bw-pir}).  A commutative ring is a
field if it models \(\forall x\,\exists y\,[x\neq0\to x\cdot y=1]\); a field
has characteristic \(p\), for a prime \(p\), if it models \(p1=0\), and
characteristic \(0\) if it models every member of the list
\((10)_p\): \(p1\neq0\), one for each prime \(p\), which replaces the
non-first-order sentence \(\forall p\,[p\text{ a prime}\to p1\neq0]\) (10).
Writing \(x^2\) for \(x\cdot x\) and \(x^{n+1}\) for \(x^n\cdot x\), a field
is algebraically closed if it models every sentence
\(\forall x_0\cdots\forall x_n\,[x_n\neq0\to\exists y\,(x_n\cdot y^n+
x_{n-1}y^{n-1}+\cdots+x_1y+x_0=0)]\); the instances that express the
intended property are those with \(n\geq1\), since the instance \(n=0\)
would require \(x_0=0\) for a nonzero \(x_0\).  In the language \(\{\in\}\)
of set theory, extensionality is
\(\forall x\,\forall y\,[\forall z\,(z\in x\leftrightarrow z\in y)\to x=y]\)
\cite[ch.~A.1, \S2, pp.~13--15]{Barwise1977FirstOrder}.
\end{hbexample}

\begin{hbtheorem}{thm:p02-w02-p02-005-bw-maximal}{proved=annex:p02-w02-p02-005-bw-maximal}{A first-order equivalent of maximality}
Let \(\mathfrak R=(R,+,\cdot,0,1)\) be a commutative ring with identity
and \(I\subseteq R\) a proper ideal.  Then \(I\) is a maximal proper ideal
if and only if \((\mathfrak R,I)\) models (9)\('\).
\end{hbtheorem}

\begin{hbtheorem}{thm:p02-w02-p02-005-bw-pir}{proved=annex:p02-w02-p02-005-bw-pir}{Being a principal ideal ring is not first-order}
There is a commutative ring \(\mathfrak R\) with identity such that
\(\mathfrak R\equiv\mathbb Z=(\mathbb Z,+,\cdot,0,1)\) and \(\mathfrak R\)
has a proper ideal that is not principal.  Consequently no set of
first-order sentences of \(\{+,\cdot,0,1\}\) has as its models exactly the
principal ideal rings.
\end{hbtheorem}

\begin{hbtheorem}{thm:p02-w02-p02-005-bw-characteristic}{proved=annex:p02-w02-p02-005-bw-characteristic}{Transfer from characteristic \(0\) (Barwise, Proposition~2.8)}
For every first-order sentence \(\psi\) of \(\{+,\cdot,0,1\}\) true in all
fields of characteristic \(0\) there is an integer \(N_\psi\) such that
\(\psi\) is true in every field of characteristic \(p\) for every prime
\(p>N_\psi\).
\end{hbtheorem}

The chapter closes its survey with properties of whole theories, each
attributed to a later chapter of the same handbook and not proved in
Chapter~A.1 \cite[ch.~A.1, \S2, p.~16]{Barwise1977FirstOrder}: the theory
of abelian groups is decidable and the theory of groups is not
(Chapters~C.1 and~C.3); the theory of algebraically closed fields of a
fixed characteristic is complete, and with the Completeness Theorem this
makes the theory of algebraically closed fields of characteristic \(0\),
and hence the field of complex numbers, decidable (Theorem~7.2 of
Chapter~C.1); by G\"odel's incompleteness theorem the ring \(\mathbb Z\)
is not decidable, so every effectively listed set of true axioms about
\(\mathbb Z\) leaves some truth unproved (Chapter~D.1).  The chapter then
directs the reader to Chapter~A.2 for basic model theory, A.3 for
ultraproducts, and A.4 for model-theoretic analogues of algebraic closure
\cite[ch.~A.1, \S2, p.~17]{Barwise1977FirstOrder}.  These are reported
claims; the handbook does not prove them in this section.

\subsection{Barwise's worked examples of the formal notions}
\label{sec:p02-w02-p02-005-bw-syntax-examples}

\begin{hbexample}{def:p02-w02-p02-005-bw-syntax-cases}{Structures, terms, atomic formulas and formulas for groups and sets (Barwise, \S3)}
For \(L=\{+,0\}\) a structure has the form
\(\mathfrak M=(M,+^{\mathfrak M},0^{\mathfrak M})\) with \(M\) nonempty,
\(+^{\mathfrak M}:M\times M\to M\) and \(0^{\mathfrak M}\in M\); one usually
writes \(G\) for \(\mathfrak M\) and drops the superscripts
\cite[ch.~A.1, \S3, p.~18]{Barwise1977FirstOrder}.  Only
\cref{def:p02-w02-p02-005-bw-structure} is used here: every such triple is
a structure for \(L\), whether or not it is a group.  Strictly, the terms
of this language are expressions such as \(+(xy)\) and \(+(0+(x0))\),
abbreviated \(x+y\) and \(0+(x+0)\).  The only constant symbol is \(0\),
so the closed terms are the expressions built from \(0\) by \(+\) alone.
Because \(L\) has no relation symbol, every atomic formula is an equation
between terms, for example \((x+y=z)\), \((x+y=y+x)\) and
\((x+y)+z=x+(y+z)\).  In \(L=\{\in\}\) every term is a variable, so the
atomic formulas are exactly \((v=w)\) and \(\in(vw)\), the latter written
\(v\in w\) \cite[ch.~A.1, \S3, pp.~18--19]{Barwise1977FirstOrder}.

The expressions \((x+y=0)\), \((\exists y\,(x+y=0))\) and
\((\forall x\,(\exists y\,(x+y=0)))\) are formulas of \(\{+,0\}\).  By
\cref{def:p02-w02-p02-005-bw-syntax}, \(\FV\) of the first is
\(\{x,y\}\), of the second \(\{x,y\}-\{y\}=\{x\}\), and of the third
\(\{x\}-\{x\}=\emptyset\).  Only the third is a sentence; it is the axiom
(3) of \cref{def:p02-w02-p02-005-bw-groups}, and only it makes sense as an
axiom.  Barwise compares the binding of a variable by a quantifier with
the definite integral \(\int_0^1(x^2+2x+1)\,dx\), whose value does not
depend on \(x\), whereas the value of the expression \(x^2+2x+1\) does
\cite[ch.~A.1, \S3, pp.~19--20]{Barwise1977FirstOrder}.
\end{hbexample}

\begin{hbexample}{def:p02-w02-p02-005-bw-polynomial}{A term as a polynomial function (Barwise, \S3)}
In the language \(\{+,\cdot,0,1\}\) of rings let \(t\) be the term written
\(x^2+2x+1\), that is \(((x\cdot x)+(x+x))+1\) under the abbreviations of
\cref{sec:p02-w02-p02-005-barwise-setup}.  For a ring \(\mathfrak R\),
\cref{def:p02-w02-p02-005-bw-satisfaction} makes \(t^{\mathfrak R}\) a
function of assignments.  Since \(x\) is the only variable of \(t\), the
value depends on \(s(x)\) alone, and if \(s(x)=a\), then
\(t^{\mathfrak R}(s)=a^2+2a+1\) computed with the operations of
\(\mathfrak R\).  So \(t^{\mathfrak R}\) is the polynomial function from
\(\mathfrak R\) into \(\mathfrak R\) given by \(t\)
\cite[ch.~A.1, \S3, p.~21]{Barwise1977FirstOrder}.  The computation uses
only that \(\mathfrak R\) is a structure for the ring language; the ring
axioms are needed to identify \(t^{\mathfrak R}\) with
\(a\mapsto(a+1)\cdot(a+1)\).
\end{hbexample}

The two facts quoted after \cref{def:p02-w02-p02-005-bw-substitution} are
the term and formula coincidence lemmas, proved in
\cref{lem:fol-coincidence} (the term case is the atomic step of its
proof).  That isomorphic structures are elementarily equivalent is
\cref{thm:p02-w02-p02-005-bw-iso-equiv}.  Barwise also records a
terminology: some authors call a finitely axiomatizable class an
\emph{elementary class}, EC, and an axiomatizable class elementary in the
wider sense, \(\mathrm{EC}_\Delta\) \cite[ch.~A.1, \S3, p.~22]{Barwise1977FirstOrder}.

\begin{hbtheorem}{thm:p02-w02-p02-005-bw-iso-equiv}{proved=annex:p02-w02-p02-005-bw-iso-equiv}{Isomorphic structures are elementarily equivalent (Barwise, \S3)}
If \(\mathfrak M\) and \(\mathfrak N\) are structures for \(L\) and
\(h:M\to N\) is an isomorphism, then for every formula \(\varphi\) and every
assignment \(s\) in \(\mathfrak M\),
\(\mathfrak M\models\varphi[s]\) iff \(\mathfrak N\models\varphi[h\circ s]\).
In particular \(\mathfrak M\equiv\mathfrak N\).
\end{hbtheorem}

\begin{hbexample}{def:p02-w02-p02-005-bw-ring-incomplete}{An incomplete theory (Barwise, \S2)}
The theory of commutative rings with identity (the axioms of
\cref{def:p02-w02-p02-005-bw-rings}) can be extended consistently both by
the field axiom \(\forall x\,\exists y\,[x\neq0\to x\cdot y=1]\) and by its
negation \cite[ch.~A.1, \S2, p.~16]{Barwise1977FirstOrder}.  What makes
this apply is the existence of two models: the ring \(\mathbb Q\)
satisfies the ring axioms and the field axiom, and the ring \(\mathbb Z\)
satisfies the ring axioms and the negation, since \(2\neq0\) and no
integer \(y\) has \(2y=1\).  By soundness
(\cref{thm:p02-w02-p02-005-bw-soundness}) neither extension proves a
contradiction, so the ring theory proves neither the field axiom nor its
negation: it is not complete.  Barwise contrasts this with the theory of
algebraically closed fields of a fixed characteristic, which he reports to
be complete (\cref{sec:p02-w02-p02-005-barwise-realm}).
\end{hbexample}

\subsection{Propositional logic inside first-order logic (Barwise 4.1--4.2)}
\label{sec:p02-w02-p02-005-bw-prop}

Barwise's \S4 sets out to show that a precise notion of proof captures
``\(\varphi\) follows from \(T\) by the laws of logic alone'' for
first-order \(\varphi\) and \(T\).  Whatever the intuitive notion is, it
implies \(T\models\varphi\); so it suffices to give valid rules and to
prove that \(T\models\varphi\) exactly when \(\varphi\) is provable from
\(T\).  The plan: propositional logic first (4.1--4.2), then a method due
in essence to Henkin for reducing first-order questions to propositional
ones (4.3--4.8), from which compactness and the L\"owenheim--Skolem
theorem follow, and then a Hilbert-style system (4.9) and a Gentzen-type
system (4.13) \cite[ch.~A.1, \S4, pp.~22--23]{Barwise1977FirstOrder}.

\begin{hbdefinition}{def:p02-w02-p02-005-bw-propformula}{Prime formulas and propositional formulas}
Let \(P\) be a set of objects called \emph{prime formulas}.  The
\emph{propositional formulas of \(P\)} form the smallest set of
expressions that contains every member of \(P\) and contains \(\neg A\),
\((A\land B)\), \((A\lor B)\) and \((A\to B)\) whenever it contains \(A\)
and \(B\).  The \emph{prime constituents} of \(A\) are the prime formulas
from which \(A\) is built.  In the application \(P\) is the set of
first-order formulas that are not propositional combinations of simpler
formulas: the atomic formulas and the formulas beginning with a quantifier
\cite[ch.~A.1, \S4, p.~23]{Barwise1977FirstOrder}.  For
\(P=\{p,q,r\}\), the expressions \(p\), \(q\), \((p\lor q)\),
\((q\lor p)\) and \(((p\lor q)\to(q\lor p))\) are propositional formulas.
\end{hbdefinition}

\begin{hbdefinition}{def:p02-w02-p02-005-bw-truth-assignment}{Truth assignments and their extension}
Let \(\mathbf t\) and \(\mathbf f\) be distinct new symbols.  A
\emph{truth assignment} for \(P\) is a function
\(v:P\to\{\mathbf t,\mathbf f\}\).  Its extension \(\bar v\) to all
propositional formulas of \(P\) is defined by recursion on length:
\(\bar v(A)=v(A)\) for prime \(A\); \(\bar v(\neg A)=\mathbf f\) if
\(\bar v(A)=\mathbf t\), and \(\mathbf t\) otherwise;
\(\bar v(A\land B)=\mathbf t\) if \(\bar v(A)=\bar v(B)=\mathbf t\), and
\(\mathbf f\) otherwise; \(\bar v(A\lor B)=\mathbf t\) if
\(\bar v(A)=\mathbf t\) or \(\bar v(B)=\mathbf t\) or both, and \(\mathbf f\)
otherwise; \(\bar v(A\to B)=\mathbf f\) if \(\bar v(A)=\mathbf t\) and
\(\bar v(B)=\mathbf f\), and \(\mathbf t\) otherwise
\cite[ch.~A.1, \S4, p.~24]{Barwise1977FirstOrder}.  The clauses are
summarized by Barwise's truth table:
\begin{center}
\begin{tabular}{cc|cccc}
\(A\) & \(B\) & \(\neg A\) & \((A\land B)\) & \((A\lor B)\) & \((A\to B)\)\\
\hline
\(\mathbf t\) & \(\mathbf t\) & \(\mathbf f\) & \(\mathbf t\) & \(\mathbf t\) & \(\mathbf t\)\\
\(\mathbf t\) & \(\mathbf f\) & \(\mathbf f\) & \(\mathbf f\) & \(\mathbf t\) & \(\mathbf f\)\\
\(\mathbf f\) & \(\mathbf t\) & \(\mathbf t\) & \(\mathbf f\) & \(\mathbf t\) & \(\mathbf t\)\\
\(\mathbf f\) & \(\mathbf f\) & \(\mathbf t\) & \(\mathbf f\) & \(\mathbf f\) & \(\mathbf t\)
\end{tabular}
\end{center}
\end{hbdefinition}

\begin{hbexample}{def:p02-w02-p02-005-bw-table-two}{A two-letter truth table (Barwise, p.~25)}
For \(((\neg(p\land\neg q)\land q)\to p)\) the table, computed row by row
from \cref{def:p02-w02-p02-005-bw-truth-assignment}, is
\begin{center}
\begin{tabular}{cc|ccccc}
\(p\) & \(q\) & \(\neg q\) & \(p\land\neg q\) & \(\neg(p\land\neg q)\) &
\(\neg(p\land\neg q)\land q\) & \((\neg(p\land\neg q)\land q)\to p\)\\
\hline
\(\mathbf t\) & \(\mathbf t\) & \(\mathbf f\) & \(\mathbf f\) & \(\mathbf t\) & \(\mathbf t\) & \(\mathbf t\)\\
\(\mathbf t\) & \(\mathbf f\) & \(\mathbf t\) & \(\mathbf t\) & \(\mathbf f\) & \(\mathbf f\) & \(\mathbf t\)\\
\(\mathbf f\) & \(\mathbf t\) & \(\mathbf f\) & \(\mathbf f\) & \(\mathbf t\) & \(\mathbf t\) & \(\mathbf f\)\\
\(\mathbf f\) & \(\mathbf f\) & \(\mathbf t\) & \(\mathbf f\) & \(\mathbf t\) & \(\mathbf f\) & \(\mathbf t\)
\end{tabular}
\end{center}
Barwise prints only the \(\mathbf f\) entries; the table above restores the
\(\mathbf t\) entries.  The formula is false exactly when \(p\) is false and
\(q\) is true \cite[ch.~A.1, \S4, pp.~24--25]{Barwise1977FirstOrder}.
\end{hbexample}

\begin{hbdefinition}{def:p02-w02-p02-005-bw-tautology}{Tautology and consistency (Barwise, Definition~4.1)}
A propositional formula \(A\) of \(P\) is a \emph{tautology} if
\(\bar v(A)=\mathbf t\) for every truth assignment
\(v:P\to\{\mathbf t,\mathbf f\}\), and \emph{consistent} if
\(\bar v(A)=\mathbf t\) for some such \(v\).  A set \(T\) of propositional
formulas is \emph{consistent} (in the sense of propositional logic) if some
truth assignment \(v\) has \(\bar v(A)=\mathbf t\) for all \(A\in T\)
\cite[ch.~A.1, \S4, Definition~4.1, p.~25; p.~26]{Barwise1977FirstOrder}.
Writing \(A\leftrightarrow B\) for \((A\to B)\land(B\to A)\), the following
are tautologies: \((A\lor\neg A)\) (excluded middle),
\(\neg(A\land\neg A)\) (contradiction),
\(\neg(A\land B)\leftrightarrow(\neg A\lor\neg B)\) and
\(\neg(A\lor B)\leftrightarrow(\neg A\land\neg B)\) (de Morgan's laws), and
\(\neg\neg A\leftrightarrow A\) (double negation)
\cite[ch.~A.1, \S4, p.~25]{Barwise1977FirstOrder}.  Each depends on at
most two formulas \(A,B\), and \(\bar v\) of the whole is determined by
\(\bar v(A)\) and \(\bar v(B)\); checking the four combinations against the
table of \cref{def:p02-w02-p02-005-bw-truth-assignment} gives
\(\mathbf t\) in every case.
\end{hbdefinition}

\begin{hbexample}{def:p02-w02-p02-005-bw-table-three}{A three-letter tautology checked twice (Barwise, pp.~25--26)}
Let \(A=((p\land q)\to r)\), \(B=(\neg r\to q)\), \(C=(p\to r)\).  The
formula \([A\land B]\to C\) has the table
\begin{center}
\begin{tabular}{ccc|cccccc|c}
\(p\) & \(q\) & \(r\) & \(p\land q\) & \(A\) & \(\neg r\) & \(B\) &
\(A\land B\) & \(C\) & \([A\land B]\to C\)\\
\hline
\(\mathbf t\)&\(\mathbf t\)&\(\mathbf t\)&\(\mathbf t\)&\(\mathbf t\)&\(\mathbf f\)&\(\mathbf t\)&\(\mathbf t\)&\(\mathbf t\)&\(\mathbf t\)\\
\(\mathbf t\)&\(\mathbf t\)&\(\mathbf f\)&\(\mathbf t\)&\(\mathbf f\)&\(\mathbf t\)&\(\mathbf t\)&\(\mathbf f\)&\(\mathbf f\)&\(\mathbf t\)\\
\(\mathbf t\)&\(\mathbf f\)&\(\mathbf t\)&\(\mathbf f\)&\(\mathbf t\)&\(\mathbf f\)&\(\mathbf t\)&\(\mathbf t\)&\(\mathbf t\)&\(\mathbf t\)\\
\(\mathbf t\)&\(\mathbf f\)&\(\mathbf f\)&\(\mathbf f\)&\(\mathbf t\)&\(\mathbf t\)&\(\mathbf f\)&\(\mathbf f\)&\(\mathbf f\)&\(\mathbf t\)\\
\(\mathbf f\)&\(\mathbf t\)&\(\mathbf t\)&\(\mathbf f\)&\(\mathbf t\)&\(\mathbf f\)&\(\mathbf t\)&\(\mathbf t\)&\(\mathbf t\)&\(\mathbf t\)\\
\(\mathbf f\)&\(\mathbf t\)&\(\mathbf f\)&\(\mathbf f\)&\(\mathbf t\)&\(\mathbf t\)&\(\mathbf t\)&\(\mathbf t\)&\(\mathbf t\)&\(\mathbf t\)\\
\(\mathbf f\)&\(\mathbf f\)&\(\mathbf t\)&\(\mathbf f\)&\(\mathbf t\)&\(\mathbf f\)&\(\mathbf t\)&\(\mathbf t\)&\(\mathbf t\)&\(\mathbf t\)\\
\(\mathbf f\)&\(\mathbf f\)&\(\mathbf f\)&\(\mathbf f\)&\(\mathbf t\)&\(\mathbf t\)&\(\mathbf f\)&\(\mathbf f\)&\(\mathbf t\)&\(\mathbf t\)
\end{tabular}
\end{center}
(Barwise again prints only the \(\mathbf f\) entries.)  The last column has
no \(\mathbf f\), so the formula is a tautology.  The shorter backward
method looks for an assignment making the formula false.  That needs
\(\bar v(A)=\bar v(B)=\mathbf t\) and \(\bar v(C)=\mathbf f\); the last
forces \(v(p)=\mathbf t\) and \(v(r)=\mathbf f\); then
\(\bar v(\neg r)=\mathbf t\), so \(\bar v(B)=\mathbf t\) forces
\(v(q)=\mathbf t\); but then \(\bar v(p\land q)=\mathbf t\) and
\(v(r)=\mathbf f\) give \(\bar v(A)=\mathbf f\), a contradiction.  So no
falsifying assignment exists \cite[ch.~A.1, \S4, pp.~25--26]{Barwise1977FirstOrder}.
In this example every step was forced by a clause of
\cref{def:p02-w02-p02-005-bw-truth-assignment}; when a step is not forced,
the backward method must try both values, and it then decides the question
only after all branches close.
\end{hbexample}

\begin{hbtheorem}{thm:p02-w02-p02-005-bw-prop-compactness}{proved=annex:p02-w02-p02-005-bw-prop-compactness}{Compactness for propositional logic (Barwise, Theorem~4.2)}
A set \(T\) of propositional formulas of \(P\) is consistent if and only
if every finite subset of \(T\) is consistent.
\end{hbtheorem}

The annex gives both of Barwise's proofs of the nontrivial half.  The
first extends \(T\) to a maximal finitely consistent set, by enumeration
when \(P\) is countable and along a well-ordering of the formulas in
general; for well-ordered \(P\) no choice is used, and the axiom of choice
supplies the well-ordering of an arbitrary \(P\).  The second applies the
Tychonoff theorem to the space \(2^P\) of truth assignments; it is faster
but hides the construction, which is the part reused later in model
theory \cite[ch.~A.1, \S4, pp.~26--28]{Barwise1977FirstOrder}.  The
statement is that of \cref{thm:prop-compactness} with an arbitrary set of
prime formulas in place of propositional variables.

\begin{hbtheorem}{thm:p02-w02-p02-005-bw-map-colouring}{proved=annex:p02-w02-p02-005-bw-map-colouring}{Map colouring and Tychonoff for \(2^P\) (Barwise, p.~28)}
\begin{enumerate}
\item Let \(k\geq1\) and let a map be given by a set \(X\) of countries
      and a symmetric adjacency relation \(E\subseteq X\times X\).  If
      every finite submap \((X_0,E\cap(X_0\times X_0))\), \(X_0\subseteq X\)
      finite, can be coloured with \(k\) colours so that adjacent
      countries get different colours, then so can the whole map.
\item For every set \(P\), the product \(2^P\) of copies of the discrete
      space \(\{\mathbf t,\mathbf f\}\) is compact.  The derivation from
      \cref{thm:p02-w02-p02-005-bw-prop-compactness} uses its first proof,
      which does not use the Tychonoff theorem.
\end{enumerate}
\end{hbtheorem}

\subsection{Henkin constants: reducing first-order logic to propositional logic (Barwise 4.3--4.8)}
\label{sec:p02-w02-p02-005-bw-henkin}

Let \(L\) be a language and let \(P\) be the set of formulas of \(L\) that
are atomic or begin with \(\forall\) or \(\exists\).  A formula of \(L\) is a
\emph{tautology of first-order logic} if it is a propositional tautology
over this \(P\).  Thus \(\forall x\,R(x)\lor\neg\forall x\,R(x)\) and
\((\forall x\,R(x)\land\exists x\,S(x))\leftrightarrow
\neg(\neg\forall x\,R(x)\lor\neg\exists x\,S(x))\) are tautologies, while
\((c=c)\) and \(\forall x\,(R(x)\lor\neg R(x))\) are prime and so not
tautologies, and \(\neg\exists x\,S(x)\to\forall x\,\neg S(x)\) has the
propositional form \(\neg p\to q\) with \(p,q\) prime and is not one
either \cite[ch.~A.1, \S4, p.~28]{Barwise1977FirstOrder}.  The three
non-examples are satisfied by every assignment in every structure, so the
tautologies capture only a small part of the laws of logic.

\begin{hblemma}{thm:p02-w02-p02-005-bw-lemma-43}{proved=annex:p02-w02-p02-005-bw-lemma-43}{Satisfaction induces a truth assignment (Barwise, Lemma~4.3)}
Let \(\mathfrak M=(M,\ldots)\) be a structure for \(L\) and \(s\) an
assignment in \(\mathfrak M\).  There is a truth assignment \(v\) to the
prime formulas of \(L\) such that for all formulas \(\varphi\) of \(L\),
\(\mathfrak M\models\varphi[s]\) iff \(\bar v(\varphi)=\mathbf t\).  In
particular any set of sentences true in \(\mathfrak M\) is consistent in
the sense of propositional logic, and every tautology of first-order logic
is satisfied by every assignment in every structure.
\end{hblemma}

\begin{hbtheorem}{thm:p02-w02-p02-005-bw-43-converse}{proved=annex:p02-w02-p02-005-bw-43-converse}{The converse of Lemma~4.3 fails (Barwise, p.~29)}
For unary relation symbols \(R,S\), the set
\(\{\forall x\,(R(x)\to S(x)),\ \forall x\,R(x),\ \exists x\,\neg S(x)\}\)
is consistent in the sense of propositional logic and has no model.
\end{hbtheorem}

All three sentences are prime, so propositional consistency sees nothing
of their quantificational structure \cite[ch.~A.1, \S4, p.~29]{Barwise1977FirstOrder}.

\begin{hbdefinition}{def:p02-w02-p02-005-bw-equality-axioms}{Equality axioms (Barwise, 4.4)}
Let \(u,w,u_1,\ldots\) range over variables and constant symbols of \(L\).
The \emph{equality axioms} of \(L\) are the formulas
\[
\begin{gathered}
(u=u),\qquad (u=w)\to(w=u),\qquad (u_1=u_2\land u_2=u_3)\to(u_1=u_3),\\
(u_1=w_1\land\cdots\land u_n=w_n)\to(R(u_1\cdots u_n)\to R(w_1\cdots w_n)),\\
(u_1=w_1\land\cdots\land u_n=w_n)\to(t(u_1\cdots u_n)=t(w_1\cdots w_n)),
\end{gathered}
\]
for every \(n\)-ary relation symbol \(R\) of \(L\) and every \(n\)-ary term
\(t\) of \(L\).  The last four might be called Leibniz Law
\cite[ch.~A.1, \S4, 4.4, p.~29]{Barwise1977FirstOrder}.  Each is satisfied
by every assignment in every structure (part 1 of
\cref{thm:p02-w02-p02-005-bw-soundness}).
\end{hbdefinition}

\begin{hbdefinition}{def:p02-w02-p02-005-bw-witnessing}{Witnessing expansion, Henkin axioms, \(T_{\mathrm{Henkin}}\) (Barwise, 4.5--4.6)}
Put \(C_0=\emptyset\) and \(L_n=L\cup C_n\).  For each formula
\(\varphi(v)\) of \(L_0\) with one free variable \(v\), let
\(c_{\varphi(v)}\) be a distinct new constant symbol, and let \(C_1\) be the
set of these.  Given \(C_n\), assign a distinct new constant symbol
\(c_{\varphi(v)}\) to each such formula \(\varphi(v)\) of \(L_n\) that is
not already a formula of \(L_{n-1}\), and let \(C_{n+1}\) be \(C_n\)
together with these.  Put \(C=\bigcup_nC_n\) and \(L(C)=L\cup C\), the
\emph{witnessing expansion} of \(L\).  The constant \(c_{\varphi(v)}\) is a
\emph{witnessing constant}; the sentences
\[
\text{I}\quad(\exists v\,\varphi(v))\to\varphi(c_{\varphi(v)}),\qquad
\text{II}\quad\varphi(c_{\neg\varphi(v)})\to\forall v\,\varphi(v)
\]
are the \emph{Henkin axioms} of types I and II.  \(T_{\mathrm{Henkin}}\)
is the set of sentences of \(L(C)\) that are Henkin axioms or of the forms
\[
\text{III}\quad\forall v\,\varphi(v)\to\varphi(t),\qquad
\text{IV}\quad\varphi(t)\to\exists v\,\varphi(v),
\]
with \(t\) a closed term of \(L(C)\); III and IV are the \emph{quantifier
axioms}.  \(\mathrm{Eq}\) is the set of equality axioms of \(L(C)\) that
are sentences, that is, contain no variables.  If \(L\subseteq L'\) and
\(\mathfrak M'=(M,F')\) is a structure for \(L'\), then
\(\mathfrak M=(M,F'\restriction L)\) is the \emph{reduct} of
\(\mathfrak M'\) to \(L\) and \(\mathfrak M'\) an \emph{expansion} of
\(\mathfrak M\) to \(L'\).  A structure \(\mathfrak M\) for \(L(C)\) is
\emph{canonical} if every element is denoted by a member of \(C\):
\(M=\{c^{\mathfrak M}\mid c\in C\}\)
\cite[ch.~A.1, \S4, pp.~29--31]{Barwise1977FirstOrder}.

\emph{Reading of ``one free variable''.}  Barwise attaches witnessing
constants to the formulas with \emph{exactly} one free variable.  Here
\(\varphi(v)\) ranges over the pairs of a formula \(\varphi\) and a
variable \(v\) with \(\FV(\varphi)\subseteq\{v\}\), so that a vacuous
quantification \(\exists v\,\varphi\) with \(v\notin\FV(\varphi)\) also
receives Henkin axioms; for it, axiom I reads
\(\exists v\,\varphi\to\varphi\) and axiom II reads
\(\varphi\to\forall v\,\varphi\).  \Cref{thm:p02-w02-p02-005-bw-vacuous}
shows that the Main Lemma below is false under the literal reading,
because the formation rules of \cref{def:p02-w02-p02-005-bw-syntax} allow
vacuous quantifiers.  The cardinality of \(C\) and every later argument of
the chapter are unaffected.
\end{hbdefinition}

The idea of the Henkin axioms: if \(\exists v\,\varphi(v)\) holds, name a
witness \(c_{\varphi(v)}\); if \(\forall v\,\varphi(v)\) fails, name a
counterexample \(c_{\neg\varphi(v)}\) \cite[ch.~A.1, \S4, p.~30]{Barwise1977FirstOrder}.
\(T_{\mathrm{Henkin}}\) is not true in every \(L(C)\)-structure, but every
\(L\)-structure expands to a model of it.

\begin{hblemma}{thm:p02-w02-p02-005-bw-lemma-47}{proved=annex:p02-w02-p02-005-bw-lemma-47}{Every structure expands to a model of \(T_{\mathrm{Henkin}}\) (Barwise, Lemma~4.7)}
Let \(\mathfrak M\) be any structure for \(L\) and \(L(C)\) the witnessing
expansion of \(L\).  There is an assignment of elements of \(M\) to the
constant symbols of \(C\) such that the resulting expansion of
\(\mathfrak M\) is a model of \(T_{\mathrm{Henkin}}\).
\end{hblemma}

\begin{hbtheorem}{thm:p02-w02-p02-005-bw-main-lemma}{proved=annex:p02-w02-p02-005-bw-main-lemma}{Main Lemma: reduction to propositional logic (Barwise, 4.8)}
Let \(L\) be a first-order language and \(L(C)\) its witnessing expansion.
For every set \(T\) of sentences of \(L\) the following are equivalent:
\begin{enumerate}
\item[(i)] \(T\) has a model, that is, some \(L\)-structure is a model of
      every sentence of \(T\);
\item[(ii)] some canonical \(L(C)\)-structure is a model of every sentence
      of \(T\);
\item[(iii)] \(T\cup T_{\mathrm{Henkin}}\cup\mathrm{Eq}\) is consistent in
      the sense of propositional logic.
\end{enumerate}
\end{hbtheorem}

The equivalence of (i) and (iii) reduces the existence of models to a
propositional question, at a price: even for finite \(T\) the
propositional theory in (iii) is infinite.  The model in
(iii)\(\Rightarrow\)(ii) is built out of the witnessing constants, as a
group is presented by generators and relations: the constants play the
generators and the sentences given the value \(\mathbf t\) play the
relations \cite[ch.~A.1, \S4, pp.~31--32]{Barwise1977FirstOrder}.
Barwise derives compactness (\cref{thm:p02-w02-p02-005-bw-compactness})
and the L\"owenheim--Skolem theorem (\cref{thm:p02-w02-p02-005-bw-ls}) from
the Main Lemma; both routes are written out in their annex proofs, beside
the derivations from \cref{thm:fol-compactness,thm:downward-ls}.  He adds
that other proofs of compactness appear in Chapters A.2 and A.3 of his
handbook and that it also follows from the Completeness Theorem below
\cite[ch.~A.1, \S4, p.~33]{Barwise1977FirstOrder}.

\begin{hbtheorem}{thm:p02-w02-p02-005-bw-vacuous}{proved=annex:p02-w02-p02-005-bw-vacuous}{The Main Lemma needs Henkin axioms for vacuous quantifiers}
Let \(L=\{d\}\) with \(d\) a constant symbol, and build \(L(C)\),
\(T_{\mathrm{Henkin}}\) and \(\mathrm{Eq}\) with witnessing constants only
for formulas with exactly one free variable.  Then
\(T=\{\exists x\,\neg(d=d)\}\) has no model, yet
\(T\cup T_{\mathrm{Henkin}}\cup\mathrm{Eq}\) is consistent in the sense of
propositional logic.
\end{hbtheorem}

\subsection{The Hilbert-style system H (Barwise 4.9--4.11)}
\label{sec:p02-w02-p02-005-bw-hilbert}

There are several distinct approaches to the Completeness Theorem, one for
each way of thinking about proofs, but only one notion of provability for
first-order logic.  A Hilbert-style system puts the weight on logical
axioms and keeps the rules to a minimum; Barwise notes that the Fitch-type
subordinate proof method and Beth's semantic tableaux (treated at length
in Smullyan's book) suit the classroom better, and that taking
\(\exists x\) as \(\neg\forall x\,\neg\) would make the system superficially
simpler \cite[ch.~A.1, \S4, p.~34]{Barwise1977FirstOrder}.

\begin{hbdefinition}{def:p02-w02-p02-005-bw-system-h}{The system H (Barwise, pp.~34--35)}
Fix a language \(L\); all formulas and terms are of \(L\), and
\(\varphi(t)\) is \(\varphi(t/v)\) under the proviso of
\cref{def:p02-w02-p02-005-bw-substitution}.  The axioms of H are
\begin{enumerate}
\item all tautologies (\cref{sec:p02-w02-p02-005-bw-henkin});
\item all equality axioms (\cref{def:p02-w02-p02-005-bw-equality-axioms});
\item all formulas \(\forall v\,\varphi(v)\to\varphi(t)\) and
      \(\varphi(t)\to\exists v\,\varphi(v)\).
\end{enumerate}
The rules of H are modus ponens, from \((\varphi\to\psi)\) and \(\varphi\)
infer \(\psi\), and the two generalization rules: if the variable \(v\) is
not free in \(\varphi\), then
\[
\frac{\varphi\to\psi(v)}{\varphi\to\forall y\,\psi(y)},\qquad
\frac{\psi(v)\to\varphi}{\exists y\,\psi(y)\to\varphi}.
\]
A \emph{proof of \(\varphi\) from} a set \(T\) of sentences is a finite
sequence \(\psi_1,\ldots,\psi_n\) with \(\psi_n=\varphi\), each member of
which is an axiom of H, a member of \(T\), or follows from earlier members
by a rule; \(T\vdash\varphi\) means that such a proof exists
\cite[ch.~A.1, \S4, pp.~34--35]{Barwise1977FirstOrder}.

\emph{Reading of \(\psi(y)\) in the generalization rules.}  Here
\(\psi(y)\) is \(\psi(y/v)\), and \(y\) is either \(v\) itself or a
variable that does not occur in \(\psi(v)\).  Without this proviso the
first rule passes from the tautology
\(\neg R(y)\to(R(v)\to\neg R(y))\) to
\(\neg R(y)\to\forall y\,(R(y)\to\neg R(y))\), which fails in a structure
with \(R^{\mathfrak M}=\{a\}\) at an assignment sending \(y\) to an element
\(b\neq a\).
\end{hbdefinition}

\begin{hbtheorem}{thm:p02-w02-p02-005-bw-completeness-h}{proved=annex:p02-w02-p02-005-bw-completeness-h}{G\"odel's Completeness Theorem for H (Barwise, 4.9)}
Let \(T\) be a set of sentences of \(L\).  A sentence \(\varphi\) is
provable from \(T\) if and only if \(\varphi\) is true in every
set-theoretic structure that is a model of \(T\): \(T\vdash\varphi\) iff
\(T\models\varphi\).
\end{hbtheorem}

\begin{hblemma}{thm:p02-w02-p02-005-bw-soundness}{proved=annex:p02-w02-p02-005-bw-soundness}{Soundness Lemma (Barwise, 4.10)}
\begin{enumerate}
\item Every axiom of H is satisfied by every assignment in every structure.
\item Let \(T\) be a set of sentences and \(\mathfrak M\) a model of \(T\).
      If \(\varphi(v_1,\ldots,v_n)\) is provable from \(T\), then
      \(\mathfrak M\models\forall v_1\cdots\forall v_n\,\varphi(v_1,\ldots,v_n)\).
\end{enumerate}
\end{hblemma}

\begin{hblemma}{thm:p02-w02-p02-005-bw-lemma-411}{proved=annex:p02-w02-p02-005-bw-lemma-411}{Derived rules of H (Barwise, 4.11, with a fifth clause)}
Let \(T\) be a set of sentences.
\begin{enumerate}
\item[(i)] If \(T\vdash(\varphi\to\psi)\) and \(T\vdash(\neg\varphi\to\psi)\),
      then \(T\vdash\psi\).
\item[(ii)] If \(T\vdash(\varphi\to\theta)\to\psi\), then
      \(T\vdash(\neg\varphi\to\psi)\) and \(T\vdash(\theta\to\psi)\).
\item[(iii)] If \(v\) does not appear in \(\psi\) and
      \(T\vdash[(\exists y\,\varphi(y)\to\varphi(v))\to\psi]\), then
      \(T\vdash\psi\).
\item[(iv)] If \(v\) does not appear in \(\psi\) and
      \(T\vdash(\varphi(v)\to\forall y\,\varphi(y))\to\psi\), then
      \(T\vdash\psi\).
\item[(v)] If \(v\) does not appear in \(\psi\) and
      \(T\vdash(\exists y\,\neg\varphi(y)\to\neg\varphi(v))\to
      ((\varphi(v)\to\forall y\,\varphi(y))\to\psi)\), then \(T\vdash\psi\).
\end{enumerate}
In (iii)--(v), \(\varphi(y)\) is \(\varphi(y/v)\) with \(y\) not occurring
in \(\varphi(v)\).
\end{hblemma}

Clauses (i)--(iv) are Barwise's \cite[ch.~A.1, \S4, 4.11, p.~35]{Barwise1977FirstOrder}.
Clause (v) is added because his proof of 4.9 removes the Henkin axioms one
at a time, on the ground that the witnessing constant of the axiom
removed first is not mentioned in the remaining ones.  That ground fails
for the pair formed by the type I axiom for \(\neg\varphi(v)\) and the type
II axiom for \(\varphi(v)\): both mention \(c_{\neg\varphi(v)}\), and both
have the same rank.  Clause (v) removes such a pair in one step; the annex
proof of \cref{thm:p02-w02-p02-005-bw-completeness-h} uses it.  Barwise
also observes that his proof uses only the derived rules (iii) and (iv),
which could therefore replace the generalization rules of H, a point
discussed by Smullyan \cite[ch.~A.1, \S4, p.~37]{Barwise1977FirstOrder}.

\subsection{The Gentzen-type systems G and G\textsuperscript{+} (Barwise 4.12--4.14)}
\label{sec:p02-w02-p02-005-bw-gentzen}

Gentzen systems put the weight on rules and reduce the logical axioms to a
minimum; they are harder to set up than a Hilbert-style system but easier
to use \cite[ch.~A.1, \S4, p.~37]{Barwise1977FirstOrder}.  The axioms
and rules of \(\mathbf G\) and \(\mathbf G^{+}\), together with the
derivation used below for Barwise's third derived claim, are collected in
the following calculus sheet; the definitions after it fix their reading.

\showcalculus{p02-w02-p02-005-rules}

\begin{hbdefinition}{def:p02-w02-p02-005-bw-sequents}{Sequents and the propositional rules (Barwise, p.~37)}
\(\Gamma\), \(\Delta\) range over finite sets of formulas.  A \emph{sequent}
is a pair \(\langle\Gamma,\Delta\rangle\), written \(\Gamma\vdash\Delta\)
and read: the conjunction of the formulas of \(\Gamma\) yields the
disjunction of the formulas of \(\Delta\); \(\Gamma,\varphi\) is
\(\Gamma\cup\{\varphi\}\).  The axioms are the sequents
\(\Gamma,\varphi\vdash\Delta,\varphi\) (rule Axiom of the calculus sheet
above), and the propositional rules are the rules \((\land\vdash)\),
\((\vdash\land)\), \((\lor\vdash)\), \((\vdash\lor)\), \((\neg\vdash)\),
\((\vdash\neg)\), \((\to\vdash)\) and \((\vdash\to)\) of that sheet.
A \emph{derivation} is a finite tree of sequents with axioms at the upper
nodes, each other sequent following from those immediately above it by a
rule \cite[ch.~A.1, \S4, p.~37]{Barwise1977FirstOrder}.  Barwise defines
derivations by a picture of a branching tree and an example rather than by
a formal definition; the description just given is the one the example
follows.  The symbol \(\vdash\) inside a sequent is distinct from
provability \(T\vdash\varphi\) in H; which is meant is fixed by whether the
left side is a finite set of formulas in a sequent or a set of sentences.
\end{hbdefinition}

\begin{hbexample}{def:p02-w02-p02-005-bw-g-derivation}{A propositional derivation (Barwise, pp.~37--38)}
A derivation of \(((\varphi\land\neg\psi)\lor\theta)\vdash(\varphi\lor\theta)\):
from the axiom \(\varphi,\neg\psi\vdash\varphi,\theta\), rule
\((\land\vdash)\) gives \((\varphi\land\neg\psi)\vdash\varphi,\theta\), and
\((\vdash\lor)\) gives \((\varphi\land\neg\psi)\vdash(\varphi\lor\theta)\);
from the axiom \(\theta\vdash\varphi,\theta\), rule \((\vdash\lor)\) gives
\(\theta\vdash(\varphi\lor\theta)\); rule \((\lor\vdash)\) applied to the
two results gives the end sequent
\cite[ch.~A.1, \S4, p.~38]{Barwise1977FirstOrder}.  Each step is an
instance of the displayed rule with the side sets \(\Gamma=\Delta=\emptyset\),
except the axioms, which are instances of
\(\Gamma,\varphi\vdash\Delta,\varphi\) with \(\Gamma=\{\neg\psi\}\),
\(\Delta=\{\theta\}\), and with \(\Gamma=\emptyset\), \(\Delta=\{\varphi\}\).
\end{hbexample}

\begin{hbtheorem}{thm:p02-w02-p02-005-bw-g-prop-completeness}{proved=annex:p02-w02-p02-005-bw-g-prop-completeness}{Completeness of the propositional fragment of G (Barwise, 4.12)}
Let \(\Gamma,\Delta\) be finite sets of propositional formulas.  The
following are equivalent:
\begin{enumerate}
\item[(i)] every truth assignment \(v\) with \(\bar v(\varphi)=\mathbf t\)
      for all \(\varphi\in\Gamma\) has \(\bar v(\psi)=\mathbf t\) for at
      least one \(\psi\in\Delta\);
\item[(ii)] \(\Gamma\vdash\Delta\) has a derivation using the axioms and
      the propositional rules.
\end{enumerate}
\end{hbtheorem}

\begin{hbdefinition}{def:p02-w02-p02-005-bw-system-g}{The first-order rules, G and G\textsuperscript{+} (Barwise, p.~39)}
To pass to first-order logic add the \emph{equality axioms}
\(\Gamma\vdash\Delta,(t=t)\) for every term \(t\); the \emph{equality
rule}, which for \(E\) either \((t_1=t_2)\) or \((t_2=t_1)\) passes
from \(\Gamma,\varphi(t_1)\vdash\Delta,\psi(t_1)\) to
\(\Gamma,E,\varphi(t_2)\vdash\Delta,\psi(t_2)\) (Equality axiom and
Equality rule of the calculus sheet above); and the quantifier rules
\((\forall\vdash)\), \((\vdash\exists)\), \((\vdash\forall)\) and
\((\exists\vdash)\) of that sheet, where in \((\vdash\forall)\) and \((\exists\vdash)\) the variable \(v\) may
not occur free in \(\Gamma\cup\Delta\), and \(\varphi(y)\) is read as in
\cref{def:p02-w02-p02-005-bw-system-h}.  The system \(\mathbf G\) has these
axioms and rules together with those of
\cref{def:p02-w02-p02-005-bw-sequents}.  The system
\(\mathbf G^{+}\) has in addition the \emph{cut} rule of the sheet, which
passes from \(\Gamma,\varphi\vdash\Delta\) and \(\Gamma\vdash\Delta,\varphi\)
to \(\Gamma\vdash\Delta\); it is the counterpart of modus ponens.
A formula \(\varphi\) is \emph{derivable from} a set \(T\) of sentences in
\(\mathbf G^{+}\) if \(\Gamma\vdash\{\varphi\}\) has a derivation for some
finite \(\Gamma\subseteq T\) \cite[ch.~A.1, \S4, p.~39]{Barwise1977FirstOrder}.
In the equality rule \(\varphi(x)\) and \(\psi(x)\) are formulas in a
variable \(x\) and the displayed instances are substitutions for \(x\).
\end{hbdefinition}

\begin{hbtheorem}{thm:p02-w02-p02-005-bw-completeness-g}{proved=annex:p02-w02-p02-005-bw-completeness-g}{Completeness Theorem for G\textsuperscript{+} (Barwise, 4.13)}
A sentence \(\varphi\) is derivable from a set \(T\) of sentences in
\(\mathbf G^{+}\) iff \(T\models\varphi\).
\end{hbtheorem}

Barwise's proof combines the Main Lemma, propositional compactness and
4.12 with five derived claims about \(\mathbf G^{+}\); the third of them
is proved by the derivation displayed as the last entry of the calculus
sheet above, which passes from \(\Gamma,(\varphi\to\theta)\vdash\Delta\)
to \(\Gamma\vdash\Delta,\varphi\)
\cite[ch.~A.1, \S4, pp.~39--40]{Barwise1977FirstOrder}, where (2) is
weakening.  In \(\mathbf G^{+}\) weakening needs a side condition on
eigenvariables, and the shared witnessing constant that calls for clause
(v) of \cref{thm:p02-w02-p02-005-bw-lemma-411} also calls for a sixth
claim here; the annex proof of \cref{thm:p02-w02-p02-005-bw-completeness-g}
states and proves both.

The cut rule is used heavily in that proof, though not in the propositional
part.  Barwise states that it can be eliminated:
\emph{Cut-Elimination Theorem (Gentzen, Barwise 4.14): any sequent that has
a derivation with the cut rule has one without it; that is,
\(\mathbf G^{+}\) and \(\mathbf G\) have the same derivable sequents}
\cite[ch.~A.1, \S4, p.~40]{Barwise1977FirstOrder}.  He indicates two
routes: a direct completeness proof for \(\mathbf G\) expanding the proof
of 4.12, as in Smullyan's book, and Gentzen's own complicated double
induction transforming derivations, whose analysis bounds the induction
provable in first-order arithmetic and is treated in his Chapter D.2 on
cut-elimination for a simpler system on formulas in negation normal form.
Neither route is carried out for \(\mathbf G\) in this section, so 4.14 is
recorded here as Barwise's claim.  The handbook's cut-elimination theorem
\cref{thm:p07-cut-elimination} concerns the calculus LK of its Part 7,
whose rules differ from those of \(\mathbf G\) (sequents of sets, and the
equality rule above).

\begin{hbexample}{def:p02-w02-p02-005-bw-exercise}{Barwise's exercise: \(\exists y\,[\varphi(y)\to\forall x\,\varphi(x)]\)}
Barwise leaves as an exercise a proof of
\(\exists y\,[\varphi(y)\to\forall x\,\varphi(x)]\) in H, G and
\(\mathbf G^{+}\), with the informal proof: choose \(y\) with
\(\neg\varphi(y)\) if there is one, and \(y\) arbitrary otherwise
\cite[ch.~A.1, \S4, p.~40]{Barwise1977FirstOrder}.  The annex item for
\cref{thm:p02-w02-p02-005-bw-exercise} writes out a derivation in
\(\mathbf G\), which is also one in \(\mathbf G^{+}\), and a proof in H, for
\(\varphi(x)\) with \(\FV(\varphi(x))\subseteq\{x\}\) and \(y\) not
occurring in \(\varphi(x)\).
\end{hbexample}

\begin{hbtheorem}{thm:p02-w02-p02-005-bw-exercise}{proved=annex:p02-w02-p02-005-bw-exercise}{The exercise is provable in G and in H}
Let \(\varphi(x)\) have \(\FV(\varphi(x))\subseteq\{x\}\), let \(y\) not occur
in \(\varphi(x)\), and write \(\varphi(y)\) for \(\varphi(y/x)\).  The
sequent \(\emptyset\vdash\exists y\,[\varphi(y)\to\forall x\,\varphi(x)]\)
has a derivation in \(\mathbf G\), and
\(\emptyset\vdash\exists y\,[\varphi(y)\to\forall x\,\varphi(x)]\) holds in H.
\end{hbtheorem}

\subsection{Beyond first-order logic (Barwise \S5)}
\label{sec:p02-w02-p02-005-bw-beyond}

Many logicians hold that there is no logic beyond first-order logic: once
all mathematical assumptions are made explicit, they can be written as
first-order axioms, and the informal notion of provable is made precise by
first-order provability.  Following a suggestion of Martin Davis, Barwise
calls this view \emph{Hilbert's Thesis}.  Its first part, that classical
mathematics is expressible in first-order logic, rests on empirical
evidence; its second part follows from the first and the Completeness
Theorem \cite[ch.~A.1, \S5, p.~41]{Barwise1977FirstOrder}.  In practice
the Thesis is not followed, as two examples show.  The torsion axiom
\(\forall x\,\exists n\geq1\,(nx=0)\) is not first-order
(\cref{thm:p02-w02-p02-005-bw-torsion}); applying the Thesis would
require axiomatizing enough number theory to carry out the arguments, so
that the theory of torsion groups would contain all of first-order number
theory.  The notions of metric space and Hilbert space are simple modulo
the ordered field \(\mathbb R\), which is categorical only relative to set
theory (\cref{thm:p02-w02-p02-005-bw-reals}); formulating everything about
metric spaces inside a formal set theory would import the problems of set
theory into a subject that is mostly first-order modulo \(\mathbb R\)
\cite[ch.~A.1, \S5, p.~41]{Barwise1977FirstOrder}.  The logics below absorb
such a notion or structure into the logic itself
\cite[ch.~A.1, \S5, p.~42]{Barwise1977FirstOrder}.

\begin{hbdefinition}{def:p02-w02-p02-005-bw-many-sorted}{Many-sorted logic and \(\mathfrak N\)-logic (Barwise, 5.1--5.2)}
Two-sorted first-order logic has two distinct sorts of variables.  A
two-sorted structure \((\mathfrak M,\mathfrak N,\ldots)\) consists of two
ordinary structures together with functions and relations on the union of
their domains.  A vector space is the two-sorted structure
\((\mathfrak F,\mathfrak V,\cdot)\): a field \(\mathfrak F\) over which the
variables \(r,s,t,\ldots\) range, the structure
\(\mathfrak V=(V,+,0)\) of vectors over which \(v,w,\ldots\) range, and
scalar multiplication \(\cdot\).  The structure
\((\mathfrak M,\mathfrak N,\ldots)\) corresponds to the one-sorted
structure \((M\cup N,M,N,\ldots)\) with unary predicates for the two sorts.
If the second structure is held fixed, equal to a given \(\mathfrak N\), the
result is \emph{\(\mathfrak N\)-logic}: its models are the two-sorted
structures \((\mathfrak M,\mathfrak N,\ldots)\) with that \(\mathfrak N\).
\(\mathbb R\)-logic suits metric spaces and real Hilbert spaces, and for the
natural numbers \(\mathbb N\), \(\mathbb N\)-logic is called
\emph{\(\omega\)-logic} \cite[ch.~A.1, \S5, 5.1--5.2, p.~42]{Barwise1977FirstOrder}.
\end{hbdefinition}

The one-sorted reduction is \cref{thm:p02-many-sorted-translation}; it
transfers most results of first-order logic to many-sorted logic and is
part of the evidence for Hilbert's Thesis.  Barwise credits Malcev and
Feferman, among others, with stressing the advantages of working directly
in the many-sorted setting, and recommends Feferman's survey of many-sorted
interpolation theorems \cite{Feferman1974ManySorted} as an introduction
\cite[ch.~A.1, \S5, p.~42]{Barwise1977FirstOrder}.

\begin{hbexample}{def:p02-w02-p02-005-bw-euclidean}{Euclidean rings in \(\omega\)-logic (Barwise, 5.2)}
A Euclidean ring is a ring \(\mathfrak R\) with a function
\(d:R\to\mathbb N\) satisfying the usual first-order axioms; so it is the
two-sorted structure \((\mathfrak R,\mathbb N,d)\), and the class of
Euclidean rings is the class of models of a set of sentences of
\(\omega\)-logic \cite[ch.~A.1, \S5, p.~42]{Barwise1977FirstOrder}.  For
example, with \(\leq\) and \(<\) of \(\mathbb N\) in the language of the
fixed sort, the division property is the two-sorted first-order sentence
\(\forall a\,\forall b\,[b\neq0\to\exists q\,\exists r\,(a=q\cdot b+r\land
(r=0\lor d(r)<d(b)))]\).  What \(\omega\)-logic supplies is that \(d\)
takes values in the true \(\mathbb N\), so that, for instance, every
descending sequence of values \(d(r_0)>d(r_1)>\cdots\) is finite, which is
what makes the Euclidean algorithm terminate.
\end{hbexample}

\begin{hbtheorem}{thm:p02-w02-p02-005-bw-n-logic}{proved=annex:p02-w02-p02-005-bw-n-logic}{Compactness fails for \(\mathfrak N\)-logic with \(\mathfrak N\) infinite (Barwise, 5.2)}
Let \(\mathfrak N\) be an infinite structure.  There is a set \(T\) of
sentences of \(\mathfrak N\)-logic every finite subset of which has a model
of \(\mathfrak N\)-logic, while \(T\) has none.
\end{hbtheorem}

\begin{hbdefinition}{def:p02-w02-p02-005-bw-weak-so}{Weak second-order logic over \(\mathrm{HF}(\mathfrak M)\) (Barwise, 5.3)}
Let \(L\) be a language with variables \(x,y,z\).  The two-sorted language
\(L^{*}\) adds set variables \(a,b,c\) and a membership symbol \(\in\).
For a structure \(\mathfrak M=(M,\ldots)\) for \(L\), put
\[
\begin{gathered}
\mathrm{HF}_0(M)=\emptyset,\qquad
\mathrm{HF}(M)=\bigcup_n\mathrm{HF}_n(M),\\
\mathrm{HF}_{n+1}(M)=\{\text{all finite subsets of }M\cup\mathrm{HF}_n(M)\},
\end{gathered}
\]
and \(\mathrm{HF}(\mathfrak M)=(\mathfrak M,\mathrm{HF}(M),\in\restriction(M\cup\mathrm{HF}(M)))\),
the \emph{structure of hereditarily finite sets on \(\mathfrak M\)}.  Weak
second-order logic (more accurately, weak finite-type logic) uses the
formulas of \(L^{*}\) with the set variables ranging over \(\mathrm{HF}(M)\)
\cite[ch.~A.1, \S5, 5.3, p.~43]{Barwise1977FirstOrder}.  The construction is
read with the elements of \(M\) as atoms: \(M\cap\mathrm{HF}(M)=\emptyset\)
and no element of \(M\) has a member in \(M\cup\mathrm{HF}(M)\).  The source
leaves this implicit; it is needed for \(\in\) to mean membership in a
hereditarily finite set.
\end{hbdefinition}

Barwise notes that the natural numbers and finite sequences are definable
in \(\mathrm{HF}(\mathfrak M)\), that \(\mathrm{HF}(\mathfrak M)\) is
admissible, so that functions can be defined by recursion (his Chapter A.7
on admissible sets, in particular 3.1 and 2.16), that the sentences of his
\S2 called weak second-order are weak second-order in this precise sense,
and that weak second-order logic has essentially the strength of
\(\omega\)-logic while being more natural in algebra
\cite[ch.~A.1, \S5, p.~43]{Barwise1977FirstOrder}.  These are recorded as
his claims; the definability and admissibility claims are not proved in
this section.

\begin{hbtheorem}{thm:p02-w02-p02-005-bw-weak-so-ls}{proved=annex:p02-w02-p02-005-bw-weak-so-ls}{L\"owenheim--Skolem holds and compactness fails for weak second-order logic (Barwise, pp.~10, 45)}
Let \(L\) be finite or countable.
\begin{enumerate}
\item If \(T\) is a countable set of sentences of \(L^{*}\) and
      \(\mathrm{HF}(\mathfrak M)\models T\) for some \(L\)-structure
      \(\mathfrak M\), then \(\mathrm{HF}(\mathfrak A)\models T\) for some
      countable \(L\)-structure \(\mathfrak A\); indeed \(\mathfrak A\) can
      be taken to be a substructure of \(\mathfrak M\) with
      \(\mathrm{HF}(\mathfrak A)\) satisfying the same \(L^{*}\)-sentences as
      \(\mathrm{HF}(\mathfrak M)\).
\item There is a set \(T\) of \(L^{*}\)-sentences, for \(L=\emptyset\), every
      finite subset of which is true in some \(\mathrm{HF}(\mathfrak M)\),
      while \(T\) is true in none.
\end{enumerate}
\end{hbtheorem}

\begin{hbcorollary}{thm:p02-w02-p02-005-bw-countable-field}{proved=annex:p02-w02-p02-005-bw-countable-field}{A countable field with the weak second-order properties of \(\mathbb R\) (Barwise, p.~12)}
There is a countable ordered field \(\mathfrak A\) such that
\(\mathrm{HF}(\mathfrak A)\) and \(\mathrm{HF}(\mathbb R)\) satisfy the same
sentences of \(L^{*}\) for \(L=\{+,\cdot,<,0,1\}\).  Among these is the
weak second-order form of the Archimedean axiom.
\end{hbcorollary}

Barwise derives the corollary from the proof of
\cref{thm:p02-w02-p02-005-bw-reals}; his sentence
\(\forall x\,\exists n\,[x\leq n1]\) is weak second-order because \(n\)
ranges over the true natural numbers \cite[ch.~A.1, \S2, p.~12]{Barwise1977FirstOrder}.
He also states at p.~10 that compactness fails for second-order and weak
second-order logic, ``as the proof of 2.1 will show'', and that the
L\"owenheim--Skolem theorem fails for second-order logic
\cite[ch.~A.1, \S2, p.~10]{Barwise1977FirstOrder}.  For full second-order
logic both failures are \cref{thm:p02-sol-full-noncompact} and the
categoricity of the second-order axioms for \(\mathbb R\),
\cref{thm:second-order-reals-categorical}, whose models are all
uncountable.

\begin{hbdefinition}{def:p02-w02-p02-005-bw-infinitary}{\(L_{\omega_1\omega}\) and the quantifiers \(Q_0\), \(Q\) (Barwise, 5.4--5.5)}
\(L_{\omega_1\omega}\) adds to the formation rules of first-order logic:
if \(\Phi\) is a countable set of formulas, then \(\bigwedge\Phi\) and
\(\bigvee\Phi\) are formulas, satisfied by \(s\) iff every, respectively
some, member of \(\Phi\) is.  Countable (\(<\omega_1\)) conjunctions and
disjunctions are allowed, but only finite (\(<\omega\)) strings of
quantifiers; an example is \(\forall x\,[x=0\lor2x=0\lor\cdots]\).
Following Mostowski, a new quantifier symbol \(Q\) may be added with the
rule: if \(\varphi(x)\) is a formula, so is \(Qx\,\varphi(x)\).  In
\(L(Q_0)\), \(\mathfrak M\models Qx\,\varphi(x)\) iff infinitely many \(x\)
satisfy \(\varphi\) in \(\mathfrak M\); in the logic Barwise writes
\(L(Q)\), with \(Q\) meaning uncountable, \(\mathfrak M\models Qx\,\varphi(x)\)
iff uncountably many \(x\) satisfy \(\varphi\)
\cite[ch.~A.1, \S5, pp.~43--45]{Barwise1977FirstOrder}.  For a formula
with further free variables the clauses are read at each assignment.
\end{hbdefinition}

Barwise reads \(L_{\omega_1\omega}\) as expressing what is first-order
modulo a countable amount of information, and remarks that all finitary
propositional connectives are already definable from \(\lor\) and \(\neg\)
\cite[ch.~A.1, \S5, p.~44]{Barwise1977FirstOrder}.

\begin{hbtheorem}{thm:p02-w02-p02-005-bw-extended-ls}{proved=annex:p02-w02-p02-005-bw-extended-ls}{L\"owenheim--Skolem and compactness for \(L_{\omega_1\omega}\), \(L(Q_0)\) and \(L(Q)\) (Barwise, pp.~44--45)}
Let the language be finite or countable.
\begin{enumerate}
\item If a countable set of sentences of \(L_{\omega_1\omega}\) has a model,
      then it has a countable model.
\item If a countable set of sentences of \(L(Q_0)\) has a model, then it has
      a countable model.
\item Compactness fails for \(L_{\omega_1\omega}\) and for \(L(Q_0)\): in each
      there is a countable set of sentences every finite subset of which has
      a model, while the set has none.
\item The L\"owenheim--Skolem theorem with \(\kappa=\aleph_0\) fails for
      \(L(Q)\): the sentence \(Qx\,(x=x)\) has a model and no countable model.
\end{enumerate}
\end{hbtheorem}

Barwise's further claims about these logics are recorded here as his:
a completeness theorem for countable theories of \(L_{\omega_1\omega}\)
needs an infinitary rule of proof; weak second-order logic translates
easily into \(L_{\omega_1\omega}\) but not conversely; in weak second-order
logic some relation is implicitly but not explicitly definable, which
Beth's theorem excludes for first-order logic and \(L_{\omega_1\omega}\);
the search for the least logic containing weak second-order logic in which
implicit definability implies explicit definability leads to the admissible
fragments \(L_A\) of \(L_{\omega_1\omega}\) of his Chapter A.7; and
\(L(Q_0)\) is essentially equivalent to \(\omega\)-logic and weak
second-order logic \cite[ch.~A.1, \S5, p.~44]{Barwise1977FirstOrder}.  For
\(L(Q)\), a countable language, he reports Keisler's Completeness Theorem
and Compactness Theorem, analogous to those of first-order logic
\cite{Keisler1970}; Keisler's basic axioms, written with ``many'' for
uncountable and ``few'' for not uncountable, are
\begin{enumerate}
\item for all \(y\) there are few \(x\) such that \(x=y\);
\item if \(\varphi(x)\to\psi(x)\) for all \(x\), and many \(x\) satisfy
      \(\varphi\), then many \(x\) satisfy \(\psi\);
\item if many \(x\) satisfy \((\varphi(x)\lor\psi(x))\), then many \(x\)
      satisfy \(\varphi\) or many \(x\) satisfy \(\psi\);
\item if there are only few \(x\) for which \(\exists y\,\varphi(x,y)\), and
      for each \(x\) only few \(y\) with \(\varphi(x,y)\), then there are only
      few \(y\) for which \(\exists x\,\varphi(x,y)\)
\end{enumerate}
\cite[ch.~A.1, \S5, pp.~44--45]{Barwise1977FirstOrder}.  Neither
Keisler's Completeness Theorem nor his Compactness Theorem is proved in
this section; the validity of the four axioms and the fact that they do not
yield ``there are many \(x\)'' are \cref{thm:p02-w02-p02-005-bw-keisler-axioms}.
Barwise adds that ``many'' and ``most'' seem almost logical notions, and
that Keisler's theorem gives a precise model of one informal concept of
``many''.

\begin{hbtheorem}{thm:p02-w02-p02-005-bw-keisler-axioms}{proved=annex:p02-w02-p02-005-bw-keisler-axioms}{Keisler's axioms are valid and do not give ``many''}
Interpret ``many \(x\) satisfy \(\varphi\)'' as \(Qx\,\varphi\) in \(L(Q)\).
Then each instance of axioms (1)--(4) holds in every structure at every
assignment, and \(Qx\,(x=x)\) is not a consequence of the set of all
instances, since that set holds in every countable structure while
\(Qx\,(x=x)\) holds in none.  Axiom (4) uses that a countable union of
countable sets is countable, an instance of the axiom of choice.
\end{hbtheorem}

\begin{hbdefinition}{def:p02-w02-p02-005-bw-abstract-logic}{Abstract logics (Barwise, 5.6)}
Abstract model theory surveys the spectrum of logics and the relations
among them.  A \emph{logic} consists of a syntax and a semantics that fit
together, in that elementary syntactic operations such as renaming symbols
can be carried out and have their intended meaning.  Barwise gives no more
precise definition here and refers to his paper \cite{Barwise1974Axioms}
\cite[ch.~A.1, \S5, 5.6, p.~45]{Barwise1977FirstOrder}.  The handbook's
precise notion is that of \cref{thm:lindstrom}.
\end{hbdefinition}

Of the logics above, weak second-order logic, \(L_{\omega_1\omega}\) and
\(L(Q_0)\) satisfy the L\"owenheim--Skolem theorem with \(\kappa=\aleph_0\)
but not the Compactness Theorem, and \(L(Q)\) satisfies the Compactness
Theorem but not the L\"owenheim--Skolem theorem
(\cref{thm:p02-w02-p02-005-bw-weak-so-ls,thm:p02-w02-p02-005-bw-extended-ls};
compactness for \(L(Q)\) is Keisler's).  Barwise explains this by
\emph{Lindstr\"om's theorem (Barwise 5.7): first-order logic is the only
logic closed under \(\land\), \(\neg\), \(\exists\) that satisfies the
Compactness Theorem and the L\"owenheim--Skolem theorem}, with a proof and
further references in Barwise's paper \cite{Barwise1974Axioms}
\cite[ch.~A.1, \S5, p.~45]{Barwise1977FirstOrder}.  The precise
formulation is \cref{thm:lindstrom}; neither 5.7 nor that theorem is
proved in this section.

\subsection{The chapter's foreword, apparatus and references}
\label{sec:p02-w02-p02-005-bw-apparatus}

Barwise writes for a mathematician who knows next to nothing about
mathematical logic, for instance one who has stalled in a later chapter
for want of the basic notions; for most readers he advises a quick reading
of \S2 and of the introductions to \S\S4 and~5
\cite[ch.~A.1, \S1, p.~6]{Barwise1977FirstOrder}.  The foreword describes
logic as attending to the language of mathematics, to the ways its objects
are defined and to the laws of reasoning about them, and traces the
subject to Leibniz's dream of a universal symbolic calculus, realized in
significant part by Boole, Frege, Peano, Russell and Whitehead, Hilbert,
Skolem, G\"odel, Tarski and their followers.  The chapter's first-page
note thanks the Alfred P.~Sloan Foundation and the National Science
Foundation for support \cite[ch.~A.1, p.~5]{Barwise1977FirstOrder}.

The attributions in the chapter that this section reproduces are:
compactness to G\"odel and Malcev (Barwise 2.4, p.~10); the reduction
method of 4.3--4.8 in essence to Henkin (p.~23); the Completeness Theorem
to G\"odel (4.9); cut-elimination to Gentzen (4.14); the name Hilbert's
Thesis to a suggestion of Martin Davis (p.~41); the advantages of
many-sorted logic to Malcev and Feferman, with Feferman's survey (p.~42);
new quantifiers to Mostowski, and the completeness theorem for ``there
exist uncountably many'' to Keisler (p.~44); and the theorem 5.7 to
Lindstr\"om, with Barwise's paper as the reference for a proof (p.~45)
\cite{Barwise1977FirstOrder}.  The chapter's referrals are: Chapters A.2,
A.3 and A.4 of the same handbook for model theory, ultraproducts and
model-theoretic algebraic closure, with the Ax--Kochen and Ersov results
discussed in two of them (p.~17); Chapter A.6, and its 1.1 in particular,
for infinitesimals (p.~13); Shoenfield's book for the axioms of ZF
\cite{Shoenfield1967} (p.~15); Chapters C.1, C.3 and D.1 for the
decidability and incompleteness claims (p.~16); Chapters A.2 and A.3 for
other proofs of compactness (p.~33); Smullyan's book for tableaux, for the
derived-rule formulation of H and for a cut-free completeness proof
\cite{Smullyan1968} (pp.~34, 37, 40); Chapter D.2 for cut-elimination
(p.~40); and Chapter A.7, in particular its 3.1 and 2.16, for admissible
sets and admissible fragments (pp.~43--44).

The chapter's reference list \cite[ch.~A.1, pp.~45--46]{Barwise1977FirstOrder}
has seven entries: Barwise's ``Axioms for abstract model theory''
\cite{Barwise1974Axioms}; Chang and Keisler's \emph{Model Theory}
\cite{ChangKeisler1973}; Enderton's \emph{A Mathematical Introduction to
Logic} \cite{Enderton1972}; Feferman's ``Applications of many-sorted
interpolation theorems'' \cite{Feferman1974ManySorted}; Keisler's ``Logic
with the quantifier `there exist uncountably many''\,'' \cite{Keisler1970};
Shoenfield's \emph{Mathematical Logic} \cite{Shoenfield1967}; and
Smullyan's \emph{First-Order Logic} \cite{Smullyan1968}.
